% ============================================================
% chapter_05_discussion.tex - Chapter 5, Discussion
%
% Purpose: theoretical interpretation of the Chapter 4 results, the
%          predicted Pearson-r bound under marginal correction, and
%          the operational + reproducibility implications.
% Sections: 5.1 Theoretical Bound on Pearson Correlation,
%           5.2 Interpreting the Optimal Aggregation Window,
%           5.3 Operational Implications,
%           5.4 Reproducibility Implications, 5.5 Paths Forward,
%           5.6 Limitations and Constraints.
% Included by main_thesis.tex.
% License: CC BY 4.0 (see ../LICENSE).
% ============================================================

\chapter{DISCUSSION}\label{chap:discussion}

The results in Chapter~\ref{chap:results} reveal a structural
pattern under marginal bias correction: the four-pillar metrics
move close to the gauge target while the temporal-skill
diagnostics stay flat. This chapter argues that the pattern is
the predicted behaviour of any marginal correction family when
scored against an external paired-time-series reference, not a
failure of the correction stages.

The four pillars (value
adjustment, distribution alignment, extreme preservation, event
detection) all move close to the gauge target, with several
metrics matching to two decimal places against the independent
Indonesian Meteorological, Climatological, and Geophysical Agency (BMKG) network.
The temporal-skill diagnostics stay flat, and they are measured
in sample against CPC-UNI rather than against the stations:
Pearson $r$ holds near $0.35$ across all three corrected stages,
Root Mean Square Error (RMSE) rises slightly, and
Nash-Sutcliffe Efficiency (NSE) drifts further from unity. Against
the independent stations the same correlation holds near $0.24$
when constructed the same way. The chapter
explains why this is structural, what it implies for users, what
the open release of the framework code enables in terms of
independent re-verification and extension, where the next
methodological frontiers lie, and where the present work is
bounded.

The argument has six parts.
Section~\ref{sec:discussion-bound} states the theoretical bound on
Pearson $r$ under marginal correction and compares the in-sample
value of $r \approx 0.35$ against CPC-UNI with the empirical band
reported for
daily Integrated Multi-satellitE Retrievals for Global Precipitation Measurement (IMERG) over the Indonesian Maritime Continent \citep{ref:ramadhan2022ts}.
Section~\ref{sec:discussion-subdaily} reports a half-hourly
empirical diagnostic, using the IMERG-L source, that quantifies
how much of the apparent $r$ ceiling is attributable to a
calendar-window mismatch between the UTC-day IMERG aggregate and
the local-day morning-observation convention used by the BMKG
validation network.
Section~\ref{sec:discussion-ops} translates the bound into
operational guidance, distinguishing applications that depend on
distributional properties from those that depend on day-specific
timing.

Section~\ref{sec:discussion-reproducibility} discusses the
implications of the parallel reproducibility contribution: an
open-source pipeline with notebook-level provenance that enables
re-verification, regional extension, and parameter tuning by
independent users.
Section~\ref{sec:discussion-paths} surveys the methodological
paths forward that fall outside the marginal correction family
and could raise the temporal skill ceiling.
Section~\ref{sec:discussion-limitations} catalogues the
limitations and constraints, preserving the not-formally-optimised
caveats on the three sensitivity parameters identified in
Chapter~\ref{chap:methods}.


% ============================================================
\section{Theoretical Bound on Pearson Correlation}\label{sec:discussion-bound}
% ============================================================

The visual asymmetry between the four-pillar pattern and the
temporal skill pattern in Chapter~\ref{chap:results} is not a
verification quirk and not a failure of the correction stages. It
is the predicted behaviour of any marginal correction family when
its outputs are scored against an independent paired-time-series
reference. The argument has two components, both established in
the earlier literature on quantile-mapping bias correction, and
both of which apply to the LSEQM+DL pipeline as constructed in
Chapter~\ref{chap:methods}: Linear Scaling (LS), Empirical
Quantile Mapping (EQM), a Generalized Pareto Distribution (GPD)
tail step, and a Convolutional Neural Network (CNN) refinement,
the last a Deep Learning (DL) architecture.

Figure~\ref{fig:bound-schematic} states the argument in three
panels, demonstrated on the archive itself rather than on an
idealised example. Panel~(a) shows the rank-preservation mechanism
over January 2021 at Stasiun Meteorologi Kalianget (Jawa): the
gauge peaks on days 4 and 6 and the satellite on day 7. Although
the month is fitted as three dekads with their own
mappings, the corrected series re-orders nothing within any dekad,
so the only four crossings fall at the dekad hand-overs while the
spread moves to the gauge ($0.69$ to $0.95$). Panel~(b) generalises
to a single dekad-of-year across the whole record at that station
(every 1 to 10 January, 2001 to 2021, $n = 195$), where one
correction mapping applies: corrected against raw is a
near-monotone rising curve, the gauge a surrounding cloud, and
$r$ stays flat, $0.368$ to $0.357$, at this station under the
archived UTC label. Panel~(c) places these beside
the in-sample values of Table~\ref{tab:headline-cpc}, a per-pixel
spatial median over the land pixels averaged across the 36 dekadal
windows.

\begin{figure}[H]
\centering
\includegraphics[width=\figFiveOneWidth\textwidth]{fig_thesis_05_bound_schematic.png}
\caption[The theoretical bound on Pearson $r$ under marginal correction]{The theoretical bound on Pearson $r$ under marginal
correction, on real data at one station and over the validation
network\label{fig:bound-schematic}}
\end{figure}

The first component is rank preservation. Monotone,
rank-preserving quantile mapping, the family to which empirical quantile mapping
with Gamma kernels and GPD tail adjustment both belong, preserves the rank order of values at each
location \citep{ref:cannon2015}. If the satellite reports a higher
value on day $i$ than on day $j$ before correction, and both days
fall in the same dekad, it will report a value no lower on day $i$
than on day $j$ after correction; across a dekad boundary a
different mapping applies, which produces the dekad-boundary
crossings visible in Figure~\ref{fig:bound-schematic}a. The
day-by-day pairing of satellite values to gauge values is fixed by
the satellite product itself; no operation in a marginal correction
pipeline can move a value from one day to its neighbour. The mapping
is monotone but not strictly so, because dry-day frequency matching
sends a substantial fraction of low-rain days to exactly zero: this
creates ties, but no reversals, so the day-ordering is preserved
without the map being one-to-one.

The CNN
refinement stage relaxes the rank-preservation property at
each pixel because the $3 \times 3$ convolution mixes information
from neighbouring pixels at the same timestep, but it does not
reorder the time series at a single pixel beyond the small
residual it learns to add on top of LSEQM. The size of that
relaxation is measurable. At the station of
Figure~\ref{fig:bound-schematic}, over the $195$ paired days of the
first January dekad shown in panel~(b) and with the raw series
sorted ascending, LSEQM reorders none of the $194$ adjacent pairs,
whereas LSEQM+DL
reorders six and introduces roughly $1$~mm of spread at equal raw
value across the $21.5\%$ of days the refinement touches. The
relaxation is therefore real but marginal, and it buys no temporal
skill: $r$ against the gauge moves from $0.368$ to $0.357$ on that
sample.

Because the residual
training target is the LSEQM-to-Climate Prediction Center (CPC) Unified Gauge-Based Analysis of Global Daily Precipitation (CPC-UNI) difference rather than
a free reconstruction, the temporal ordering at each pixel after
the CNN refinement remains a small perturbation of the LSEQM
ordering. The rank-preservation argument therefore carries through
to the full pipeline at the level of pixel-wise paired time series,
which is the level at which Pearson $r$ against an external gauge
is computed.

The second component is variance inflation without skill gain.
A pointed critique of quantile mapping in regional climate
applications \citep{ref:maraun2013} established that quantile
mapping widens spread to match the marginal target without
improving time-dependent skill, because spread alignment and
day-by-day pairing are independent properties of a paired
series. Within a single dekad at a pixel, Linear Scaling is a
positive rescaling, which leaves the Pearson correlation exactly
where it was; the nonlinear EQM and GPD stages introduce no rank
reversals, so they cannot re-pair a satellite day with a different
gauge day. The exact invariance under Linear Scaling holds only at
that granularity.

A correlation pooled across dekads and pixels
mixes samples that each carry a different scaling factor, so the
pooled value shifts slightly even under Linear Scaling: against the
BMKG stations it moves by $+0.003$
(Section~\ref{sec:results-intercomparison}).
Neither step fixes the Pearson correlation exactly, because the map
is nonlinear and dry-day matching creates ties, but in practice the
correlation is observed to move only within single-pixel sampling
uncertainty. Even a perfect spread match therefore has little
purchase on $r$. The combination of rank
preservation and spread-without-skill yields the headline
proposition of this chapter:

% Keep the indented proposition together on one page: ~8 lines of
% quote text + small breathing room. If less than that remains on the
% current page, force a break so the whole block lives on the next.
\needspace{10\baselineskip}
\begin{quote}
\emph{A monotone, rank-preserving marginal correction cannot
re-order a satellite series in time, so it cannot manufacture
day-by-day agreement with an external paired gauge that the raw
series does not already have. The Pearson correlation is therefore
held near its raw value; the residual movement comes only from the
nonlinear reshaping of amplitudes and the small rank-changing
contribution of the spatial CNN refinement, and is observed to
stay within single-pixel sampling uncertainty.}
\end{quote}

\noindent The proposition is what the four-pillar pattern and the
flat $r$ pattern are jointly expressing.

The bound is not a claim that $r \approx 0.35$ is the value a daily
product must take. Reported daily correlations vary widely with the
setting: against the dense Chinese gauge analysis the median for
IMERG is near $0.75$, while over the Indonesian Maritime Continent
the daily value is lower and rises substantially only at dekadal or
monthly aggregation (Section~\ref{sec:lit-bounds}). The corrected
LSEQM+DL value of $0.348$, measured per pixel against the
in-sample CPC-UNI target at the native UTC window and averaged over
the $36$ dekadal windows, sits at the low end of that spread. What
the bound asserts is narrower and is what the stage comparison
tests: whatever the starting value, marginal correction does not
raise it.

The ceiling has two
contributors. The first is a calendar-window component: the
UTC-dated IMERG-L archive coincides with the CPC-UNI calibration
label but is offset by about one day from the local-day BMKG
morning-observation label. Against BMKG the daily correlation
is therefore suppressed under UTC alignment and recovered by re-aggregating
the satellite to the local-observation window: a fixable metadata
artefact rather than a satellite limit, quantified in
Section~\ref{sec:discussion-subdaily}. The second is the residual
timing accuracy of the underlying retrieval, which lies upstream
of every step in the pipeline and requires methods outside the
marginal correction family
(Section~\ref{sec:discussion-paths}).
Section~\ref{sec:discussion-ops} returns to the operational
significance of this split.

The pattern is
also stable across the calendar year.
Figure~\ref{fig:taylor-monthly} aggregates the BMKG Taylor
analysis by month and shows that the angular position of the
three method symbols, which encodes correlation, sits at the
same place in every month. The radial position, which
encodes the standard deviation ratio, sits near unity for LSEQM
and LSEQM+DL and short for LS. The bound is therefore not a
seasonal artefact and does not depend on the wet- or dry-season
sample composition.

\begin{figure}[H]
\centering
\includegraphics[width=\figFiveTwoWidth\textwidth]{fig_thesis_05_taylor_monthly.png}
\caption[Monthly Taylor diagram aggregated across the BMKG validation stations]{Monthly Taylor diagram aggregated across the $171$ BMKG
stations carrying a valid per-dekad correlation\label{fig:taylor-monthly}}
\end{figure}

A complementary reading uses the canonical decomposition of the
mean squared error into a bias term, a spread-mismatch term, and a
correlation-residual term \citep{ref:gupta2009}. The bias term is
moved cleanly by the Linear Scaling stage, and the spread-mismatch
term is driven toward zero by the EQM stage as the
standard-deviation ratio moves to unity and beyond. The
correlation-residual
term, equal to
$2\,\sigma_{\text{sat}}\,\sigma_{\text{gauge}}\,(1-r)$, is
the mechanism through which the RMSE actually increases between
LS ($13.1$ mm per day) and LSEQM+DL ($14.1$ mm per day). Once the
satellite spread is inflated to match the reference spread without
any improvement in $r$, this term grows, because
$\sigma_{\text{sat}}$ has risen toward $\sigma_{\text{ref}}$ while
$(1-r)$ stays near $0.65$, and its growth outweighs the vanishing
spread-mismatch term $(\sigma_{\text{sat}} - \sigma_{\text{ref}})^2$.
The whole decomposition is measured in sample against CPC-UNI,
where the standard-deviation ratio does not stop at unity but
overshoots to $1.15$; the convergence to $1.00$ reported in
Chapter~\ref{chap:results} is the out-of-sample result against the
BMKG stations, and mixing the two would misstate the mechanism.

This is the
variance-inflation-without-skill effect identified for quantile
mapping in regional climate applications \citep{ref:maraun2013}:
matching the spread without matching the timing increases the
squared error at the corrected stage relative to the
under-spread raw stage. NSE follows the same logic by
construction. The fact that three temporal-skill diagnostics
fail to improve and one improves only mildly does not indicate
failure; it indicates that the pipeline is doing exactly what
marginal correction can do, and not what it cannot.

The framing of this pattern matters for the wider field. In the
absence of stage-resolved attribution, a study that reported only
the aggregate $r$ would conclude that the LSEQM+DL framework
delivers little improvement over raw IMERG-L, and the substantial
movement in the value, distribution, extreme, and detection
pillars would be invisible. The four-pillar verification framework
of Section~\ref{sec:methods-metrics}, paired with Pearson $r$ as a
separate temporal-skill diagnostic rather than as an
all-encompassing skill score, lets the structural limit and the
methodological improvement appear in the same report. This is the
verification posture the present work argues for, and it is a
posture that the parallel reproducibility contribution
(Section~\ref{sec:discussion-reproducibility}) lets any reader
independently reconstruct on their own data.


% ============================================================
\section{Interpreting the Optimal Aggregation Window}\label{sec:discussion-subdaily}
% ============================================================

\textbf{The calendar-window mismatch.} The bound in
Section~\ref{sec:discussion-bound} predicts that
marginal correction cannot lift Pearson~$r$ above the timing
quality of the underlying retrieval, but it does not decompose
the absolute ceiling into its components. One component the bound
discussion does not isolate is the calendar-window definition used
to label each daily value. IMERG-L distributes daily totals
aggregated over the Coordinated Universal Time (UTC) calendar day~\citep{ref:imergdl07}.
The BMKG climatological-station daily totals are reported
under the morning-observation convention specified by Head of BMKG
Decree SK.32/TL.202/KB/BMG-2006 \citep{ref:bmkg2006sk32},
consistent with the World Meteorological Organization principle of
measuring daily precipitation at fixed observation times common to
the network \citep{ref:wmo2018no8}. Under that convention the
rainfall total assigned to
date~$d$ is the 24-hour accumulation ending at the BMKG morning
observation on date~$d$, nominally $07{:}00$ local time. Because
Indonesia spans the three local-time zones WIB, WITA, and WIT
(UTC$+7$, $+8$, $+9$), the BMKG daily window ends near $00{:}00$
UTC for WIB stations and a few hours earlier for stations east of
the WITA boundary. The daily IMERG-L window and the daily BMKG
window are therefore offset by approximately 24~hours in the
worst case, and the daily product values labelled by the same
calendar date refer to different 24-hour rainfall periods.

Section~\ref{sec:results-window} established empirically that
re-aggregating the half-hourly source recovers the GPM-era
correlation from $r = 0.20$ at the UTC-day window to $r = 0.57$ at
$h = -23$~h. That offset is not a fitted parameter but the
signature of the gauge convention. Under the morning-observation
rule, the total assigned to date~$d$ is the 24-hour accumulation
ending at the $07{:}00$ observation; for a WIB station this window
is $[00{:}00\,\text{UTC}\ (d{-}1),\ 00{:}00\,\text{UTC}\ (d)]$,
which begins $24$~h before the UTC day the archive labels with
date~$d$. The empirical peak at $-23$~h sits within the two-to-four
hour spread expected from varying station observation times,
retrieval blending across the half-hourly timestamp, and the
diurnal phase of maritime-continent convection.

\textbf{Convention dependence of the gauge references.} The two
gauge references differ in both the hour at which a daily total
closes and the date it is then given, and resolving the difference
sharpens the picture. The
gridded CPC-UNI archive dates each daily total by a per-cell
end-of-day field; for every Indonesian cell this field is
$24$~hours, so the total closes at $00{:}00$ UTC and is assigned
the UTC calendar date \citep{ref:cpcuni_readme}, whereas the BMKG
archive keeps the local morning-read date one calendar day later.
The CPC-UNI total therefore coincides with the UTC-day IMERG-L
window, while the BMKG label sits one calendar day from it. A day-shift
test confirms the split directly: the IMERG-versus-CPC-UNI daily
correlation is highest at the native same-date pairing and falls
when CPC-UNI is shifted by a day, whereas the IMERG-versus-BMKG
correlation is highest only after the satellite is re-aggregated
to the local-observation window (Figure~\ref{fig:convention}a).
The calendar-window offset is thus specific to the BMKG
labelling: IMERG and CPC-UNI share the UTC day and carry no
window offset, so the calibration target is already consistent
with the satellite.

At
the station locations both references reach the same
retrieval-limited correlation (approximately $0.57$ in the GPM
era), but at offset windows: CPC-UNI at the UTC day and BMKG at
$h = -23$~h. The practical consequence is that daily
satellite-gauge correlation is convention-dependent: no single
calendar label maximises agreement with both references at once,
so a corrected product is best distributed on an explicit
local-observation day with all references harmonised to it
(Figure~\ref{fig:convention}b). The
agreement attainable after harmonisation remains bounded by the
retrieval timing accuracy, so the residual ceiling cannot be
raised by relabelling and requires sub-daily or multivariate
methods (Section~\ref{sec:discussion-paths}).

\begin{figure}[H]
\centering
\includegraphics[width=\figFiveConventionWidth\textwidth]{fig_thesis_05_convention.png}
\caption[The calendar-window convention conflict and its resolution]{Daily
IMERG correlation against CPC-UNI and BMKG versus aggregation-window offset at the
$178$ station locations with valid half-hourly coverage (GPM era): (a) the convention
conflict and (b) its resolution by relabelling CPC-UNI to the local-observation day.\label{fig:convention}}
\end{figure}

The diagnostic pools the $178$ stations with valid half-hourly
coverage, a slightly larger set than the $172$ out-of-sample
validation stations used for the headline results.
Table~\ref{tab:convention} lists the four correlations directly: each
reference reaches approximately $0.57$ at its own optimal window
(CPC-UNI at the UTC day, BMKG at $h = -23$~h) and falls to $0.18$
to $0.20$ at the other, a mirror-image structure across the two
calendar labels.

\begin{table}[H]
\centering
\caption[Convention-dependence of the daily satellite-gauge correlation]{Daily
IMERG Pearson $r$ at the $178$ BMKG station locations (GPM era), under the UTC
and the local-day ($h=-23$~h) aggregation windows\label{tab:convention}}
\begin{tabular}{lcc}
\toprule
Reference & UTC window ($h=0$) & Local-day window ($h=-23$~h) \\
\midrule
CPC-UNI (calibration) & $\approx 0.57$ & $\approx 0.18$ \\
BMKG (validation)     & $\approx 0.20$ & $\approx 0.57$ \\
\bottomrule
\end{tabular}
\end{table}

\textbf{Whole-domain confirmation against CPC-UNI.} The
station-based day-shift test extends to a continuous map by
repeating the offset sweep against the gridded CPC-UNI analysis at
every $0.5^\circ$ cell (Figure~\ref{fig:gridded-cpc}). At
CPC-UNI's native UTC labels the optimal offset is $h^\star
\approx 0$ across the gauge-constrained western and central cells
(Figure~\ref{fig:gridded-cpc}a), the whole-domain signature of
CPC-UNI sharing the IMERG-L UTC day; relabelling CPC-UNI to the
local-observation day moves the optimum to $h^\star \approx -23$~h
everywhere (Figure~\ref{fig:gridded-cpc}b). The offset is
identical in the timing-imprecise TRMM-input era ($2001$ to
$2013$) and the GPM era, which is consistent with a labelling
effect rather than a property of retrieval skill. Pooled over the
$134$ cells that hold
BMKG stations the correlation reaches $0.56$ in the GPM era
(Figure~\ref{fig:gridded-cpc}c, bold), close to the
per-station diagnostic and Table~\ref{tab:convention}; pooled over
the whole domain it falls to $0.334$ (faint), close to the
in-sample whole-coverage correlation against CPC-UNI
(Table~\ref{tab:cpc-corr-dist}), diluted by the data-sparse east
where the contributing-gauge network thins and $h^\star$ turns noisy. The era gain against CPC-UNI is modest
($0.51$ to $0.56$) next to the point-gauge recovery ($0.36$ to
$0.57$) because the smoothed $0.5^\circ$ analysis forgives the
point-representativeness noise that depresses the early-era
point-gauge correlation, so both references converge near $0.56$
to $0.57$ once the retrieval resolves the diurnal cycle.

\begin{figure}[H]
\centering
\includegraphics[width=\figFiveGriddedCpcWidth\textwidth]{fig_thesis_05_gridded_cpc.png}
\caption[Gridded window offset against CPC-UNI over the full record]{Gridded
window-offset diagnostic against CPC-UNI over the full record\label{fig:gridded-cpc}}
\end{figure}

\textbf{Timezone bands and the single peak.} The single broad
peak shared across the WIB, WITA, and WIT bands,
rather than three distinct peaks at $-7$, $-8$, and $-9$~h, shows
that the morning-observation convention dominates over the
local-midnight timezone offset. The reason is that the observation
time is a legal-clock time, a step function of the timezone region,
not a solar time that varies continuously with longitude. Every
station within a band therefore shares the same UTC window
boundary, which cancels the continuous solar drift and leaves no
systematic longitudinal gradient in $h^{\star}$
(Section~\ref{sec:results-window}); the residual two-hour spread
reflects the diurnal-phase variation and the breadth of the
correlation peak. One consequence is a small artefact at timezone
boundaries: two stations a few kilometres apart but on opposite
sides of the WIB/WITA line use legal reference times one hour
apart, so their gauge windows differ by one hour in UTC despite
near-identical rainfall and near-identical solar time. This is an
intrinsic limitation of attributing a legal-clock observation to a
spatially continuous field, and it argues for defining any
operational window from the gauge reporting convention rather than
from solar geometry.

\textbf{Confinement to the GPM era.} The recovery is confined to
the GPM era, and the reason is
instructive (Figure~\ref{fig:window-precision}). Before the full
GPM constellation, the half-hourly field places a given day's rain
with roughly one-day uncertainty: its correlation with the gauge on
the day before a wet day is nearly as high as on the wet day
itself, so no realignment of the 24-hour window can recover the
day-by-day match, and the correlation stays near $r = 0.36$ at
every offset. The GPM-era field is day-precise, so aligning the
window to the gauge convention recovers the correlation. The
window-offset search therefore doubles as a diagnostic of a
product's day-to-day temporal precision, and it places the IMERG
precision improvement at the $2014$ to $2015$ transition to the
GPM constellation. A telling detail is that in the early era the
correlation optimum sits near the UTC day rather than at the
convention offset, consistent with a field too temporally smeared
to register the gauge window; the reason the optimum takes this
sign in the early era is not established here. The era dependence is
independently corroborated over Indonesia: an evaluation in
mountainous Sumatra attributes the improved daily IMERG accuracy in
the GPM era to the denser passive-microwave constellation and the
algorithm advances introduced with the 2014 GPM Core launch
\citep{ref:ramadhan2024sumatra}.

\begin{figure}[H]
\centering
\includegraphics[width=\figFiveWindowPrecisionWidth\textwidth]{fig_thesis_05_window_precision.png}
\caption[Why the window recovery is confined to the GPM era]{Day-offset
correlation (a) and window-offset correlation (b) for the pre-GPM and
GPM eras, showing why the window recovery is confined to the GPM
era.\label{fig:window-precision}}
\end{figure}

\textbf{Calibration and validation.} The two gauge references are
assembled under different conventions. CPC-UNI closes each daily
total at $00{:}00$~UTC and assigns it the UTC date
\citep{ref:cpcuni_readme}, while BMKG closes at the morning
observation and assigns the local date under
SK.32/TL.202/KB/BMG-2006 \citep{ref:bmkg2006sk32}. The two archives
therefore label a day about one day apart, and the
window mismatch surfaces at only one end of the
pipeline. The CPC-UNI label coincides with the UTC-day IMERG-L
archive, so the calibration pairing carries no window offset; the
BMKG label sits one day from the archive, so the $-23$~h offset
appears only in the validation pairing. The diagnostic was
performed against BMKG because it provides the independent
out-of-sample station network at point-pixel resolution. The
four-pillar metrics (value, distribution, extreme, detection) are
functions of the marginal distribution and of categorical events
and are therefore window-invariant when the daily sample
composition is preserved, so the calibration of LS, EQM, and GPD
against CPC-UNI remains correct under the present window choice.
Pearson~$r$ is the only metric the window choice affects, and the
effect is confined to the comparison against the local-day BMKG
gauges.

\textbf{The monthly scale and the headline ceiling.} Because the
four-pillar metrics are essentially window-invariant (the daily
sample composition being nearly preserved across offsets), only the
absolute Pearson~$r$ depends on the aggregation window, and
aggregating to the monthly scale dissolves even that dependence:
Section~\ref{sec:results-window} reports a monthly correlation near
$0.80$ whether the daily windows are UTC-aligned or gauge-matched.
The product therefore carries the monthly rainfall faithfully while
mis-assigning it across day boundaries at the daily scale, which
locates the daily ceiling in the timing of the day-by-day pairing
rather than in the amount of rainfall captured. The bound of
Section~\ref{sec:discussion-bound} separates accordingly into a
calendar-window component, the larger and recoverable part, and a
residual retrieval-skill component that no window choice removes.
The correlation against CPC-UNI, the in-sample calibration target,
holds near $0.35$ (Table~\ref{tab:cpc-corr-dist}) and is left
pinned close to its raw value by the marginal-correction pipeline, which
introduces no rank reversals in its LSEQM stages and only a small
rank-changing CNN contribution, so Pearson $r$ against CPC-UNI
changes only within sampling noise. The independent correlation
against the BMKG point gauges is lower and era-dependent; over the
GPM era it rises from $0.20$ at the UTC-day window to $0.57$ at the
gauge-matched window (Section~\ref{sec:results-window}), the
recoverable calendar-window component in action.

\textbf{Scope and bounds of the contribution.} The window-offset
search adopted here is the gauge-timing inference
of \citet{ref:beck2019}, developed for global precipitation
merging. Evaluations over the Maritime Continent pair the two
series in Coordinated Universal Time rather than on the gauge
reporting day \citep{ref:ramadhan2022ts}, and daily
bias-correction workflows likewise report correlation
against a fixed window \citep{ref:ramadhan2024sumatra, ref:suhartanto2025}. The contribution here is therefore not
the offset search but its regional grounding in the BMKG
convention, its propagation into the verification of a correction
pipeline, and the quantified recovery with its seasonal, spatial,
and era structure. A broader implication, which this single-region
result does not establish, is that daily satellite-gauge
correlations reported against a fixed UTC window may understate a
product's true day-by-day skill \citep{ref:tang2020,
ref:ramadhan2022ts}; this is offered as a hypothesis for
multi-region testing, not a conclusion.

The recovery is bounded, and the bounds govern how it should be
read. It is confined to the GPM era; it is partial, lifting the
pooled correlation to approximately $0.57$ rather than to unity; it
is an upstream re-aggregation of the raw half-hourly archive,
outside the LSEQM+DL pipeline, which by construction cannot move
$r$. It requires access to the half-hourly source, since an
already-aggregated daily archive cannot be re-windowed; the
$-23$~h offset is specific to the local-day morning convention and
must be re-derived for any other reporting practice; and it acts on
the daily correlation alone, leaving the four pillars and the
monthly skill unchanged. The operational re-aggregation of the
daily archive over a gauge-matched window is a single-stage change
to the pre-processing, identified as the most actionable path
forward in Section~\ref{sec:discussion-paths}.


% ============================================================
\section{Operational Implications}\label{sec:discussion-ops}
% ============================================================

The split between four-pillar improvement and flat temporal skill
translates directly into operational guidance for users of the
corrected IMERG-L product. The pipeline corrects what marginal
correction can correct; users should match their application to
the metrics that moved, not to the ones that did not. The
guidance organises around the question of whether the application
depends on the daily distribution or on the day-by-day pairing of
satellite values to a specific calendar day.
Table~\ref{tab:app-classes} summarises the resulting split into
two application classes, with example applications, the
dominant skill dimension each one needs, and the
relevant headline metric from Table~\ref{tab:headline}.

\begin{table}[H]
\centering
\caption[Application classes for corrected IMERG-L over Indonesia]{Application classes for corrected IMERG-L over
Indonesia\label{tab:app-classes}}
\begingroup\small
\begin{tabular}{@{}p{2cm}p{5.5cm}p{2.5cm}p{2.5cm}@{}}
\toprule
Class & Example applications & Dominant
skill dimension & Relevant headline metrics \\
\midrule
Well-served &
SPI/SPEI drought monitoring, flood-frequency analysis,
threshold-exceedance climatology, hydrological model calibration
against multi-year statistics, climate-risk insurance pricing,
agricultural dekadal water balance &
Distributional &
SDR, KS $p$, $Q_{95}$/$Q_{99}$ ratios, wet-day frequency \\
\addlinespace
Poorly-served &
Real-time flood-event nowcasting, single-event forecasting,
day-1 hydrological model initialisation, day-specific field
operations &
Day-by-day timing &
Pearson $r$, RMSE, NSE (set by the raw retrieval's timing; not improved by correction) \\
\bottomrule
\end{tabular}
\endgroup
\end{table}

The well-served class includes water-budget calculations at
dekadal-to-monthly aggregation, drought indicator computation
including the Standardized Precipitation Index (SPI) and the
Standardized Precipitation Evapotranspiration Index (SPEI),
extreme-event frequency analysis, threshold-exceedance
climatology, hydrological model calibration against multi-year
statistics, and climate risk assessment for agricultural and
infrastructure planning. All depend on distributional properties:
the wet-day frequency, the daily intensity distribution, the
upper percentiles, and the multi-year statistics that aggregate
across many days. The four-pillar results in
Chapter~\ref{chap:results} establish that the corrected product
reproduces all of these properties close to the gauge target,
typically to within $1\%$ to $5\%$ at the headline metrics.

The threshold-stratified detection result of
Section~\ref{sec:results-threshold} refines this guidance in a way
the aggregate scores conceal, and it is the one result in this
thesis with an immediate operational consequence. Aggregate
probability of detection falls under the full pipeline, which read
alone would argue against using the corrected product for any
detection task. The stratified view reverses that reading at and
above $20$~mm/day, and the reversal has a mechanism rather than a
coincidence: the aggregate is dominated by light rain, where the
uncorrected over-detection of drizzle inflates the hit count, while
the corrected product recovers the truncated upper tail that carries
the events an operational threshold is set to catch. A user whose
decision threshold sits in the moderate-to-heavy range should
therefore read the stratified curve and not the aggregate; a user
screening for any measurable rain should not. This is the same
lesson as the four-pillar split, applied within a single pillar:
an aggregate score can conceal the regime that matters.

The poorly-served class includes single-event nowcasting,
real-time flood-event prediction at a specific date, day-1
hydrological model initialisation, and any application that
requires the corrected product to identify the precise calendar
day of a heavy rain event. These applications depend on
day-by-day timing that the corrected product inherits from raw
IMERG-L, pinned close to its raw value. The Pearson $r$ ceiling is
the operational expression of this limit: $0.35$ in sample against
CPC-UNI as a per-pixel dekadal median, and, pooled over all
station-days of the GPM era, $0.20$ against the independent BMKG
stations under the archived date label, rising to $0.57$ once the
labels are matched but no further under correction. On a randomly chosen day
the corrected product carries the same timing information as the
raw retrieval, and no marginal correction can add the timing
information that the satellite did not measure.

The split tracks day-specific timing, not topic. Flood-related
applications fall on both sides of the line: flood-frequency
analysis from a $25$-year corrected record, flood-threshold
climatology against the $99$th percentile of the gauge-aligned
daily distribution, and flood-risk insurance pricing all depend
on the daily distribution and are well-served. Real-time
flood-event prediction at a specific date, in contrast, depends
on day-by-day timing and is not.

Hydrological applications split
similarly. Multi-year model calibration against gauge-aligned
statistics is well-served; day-1 forecast initialisation is not.
Agricultural applications split on the same axis. Dekadal water
balance and SPI-based drought monitoring are well-served;
single-day field operation scheduling based on a predicted heavy
rain event is not.

The corrective work at the daily scale is not optional for the
well-served class. Dekadal, monthly, and seasonal aggregates
inherit the bias structure of the daily fields they sum over. A
daily product with a $20\%$ wet-day over-count produces a
dekadal aggregate with the same fractional over-count in
wet-day counts and a similarly biased dekadal mean intensity.
The pipeline corrects the daily fields so that the downstream
aggregates derived from them inherit a bias-free input. The
operational guidance is not to skip daily correction and aggregate
the raw product, but to use the corrected daily product for
applications matched to its design and to recognise the daily
timing ceiling as a structural feature inherited from the
satellite retrieval, not a methodological gap to be closed by a
better correction.

A specific caveat applies in the Maritime Continent. Because the
UTC-dated IMERG-L archive is offset by about a day from the
local-day BMKG validation gauges \citep{ref:bmkg2006sk32}
(Section~\ref{sec:discussion-subdaily}), the Pearson $r$ ceiling
reported in Chapter~\ref{chap:results} against those gauges mixes
the underlying retrieval timing accuracy with a calendar-window
label artefact. The implication for operational users is that the
low daily correlation is largely a fixable metadata effect: the
recovery quantified in Section~\ref{sec:discussion-subdaily} is
obtained by re-aggregating the satellite to the local-observation
day, an input-side change that the CPC-UNI calibration label does
not need because it already coincides with the UTC day. The
residual at the matched window is the genuine retrieval limit.
Section~\ref{sec:discussion-paths} returns to this point as the
most operationally actionable next step for users in the UTC+7
to UTC+9 band.


% ============================================================
\section{Reproducibility Implications}\label{sec:discussion-reproducibility}
% ============================================================

The scientific result of Chapter~\ref{chap:results} and the
theoretical bound of Section~\ref{sec:discussion-bound} are
methodologically interesting only if a second reader can
reconstruct the same numbers from the same inputs and verify the
claim independently. The parallel methodological contribution of
the present work is a reproducible end-to-end pipeline designed so
that the verification posture argued for in
Section~\ref{sec:discussion-bound} is not a one-off result of the
present analysis but a procedure that any reader can apply to
their own region, their own satellite product, or their own
question. This section discusses what the open release enables, in
the categories that the literature on reproducible computational
research identifies as the substantive contributions of public
code and notebook artefacts \citep{ref:peng2011, ref:stodden2014}.

The first category is independent re-verification. The
four-pillar headline numbers in Table~\ref{tab:headline} are
generated by a specific sequence of notebook cells from a specific
set of input files. An independent reader can re-run the
notebooks against the same inputs and confirm that the reported
numbers come out of the pipeline, not out of selective reporting
or post-hoc parameter tuning. The random seeds, dependency
versions, and dekad-by-dekad runtimes needed for an independent
reconstruction of the headline table to the third decimal place
are recorded with the source release. The notebook-level
provenance,
combined with the dual local-and-Colab execution model described
in Section~\ref{sec:methods-reproducibility}, lets the
verification proceed on commodity hardware without the
institutional access barriers that the FAIR (Findable, Accessible,
Interoperable, Reusable) data principles
\citep{ref:wilkinson2016} were formulated to remove. The
verification budget for one dekad is on the order of minutes on
the Colab free tier, which makes spot-checking a feasible
classroom exercise rather than a major undertaking.

The second category is regional and product extension. The
pipeline does not encode the choice of CPC-UNI as a reference, the
choice of Indonesia as a domain, or the choice of $172$ BMKG
stations as a validation set. All three are configurable through
the \texttt{config.yml} file described in
Section~\ref{sec:methods-reproducibility}. A user in West Africa
who wishes to apply the same framework to IMERG-L over Senegal
against a local gauge network can do so by changing input paths
and domain boundaries in the configuration, without modifying the
source code. A user who wishes to substitute the Climate Hazards
Group InfraRed Precipitation with Stations (CHIRPS)
\citep{ref:funk2015} or the PERSIANN Climate Data Record
(PERSIANN-CDR) \citep{ref:ashouri2015} as the reference, or the
Multi-Source Weighted-Ensemble Precipitation (MSWEP)
\citep{ref:beck2017} as the satellite
input, can do so by changing the precipitation-variable names in
the configuration.

The framework does not guarantee that the
resulting bias correction will be skilful in those settings. The
sensitivity of the GPD shape parameter to small samples, the
calibration of $\alpha$ for blending, and the appropriateness of
the saturation count for the local gauge network all require
local revisiting. What the open release does enable is that the
revisiting can be done by domain scientists rather than by the
original authors, and that the verification of any resulting
claim is reproducible by a third party.

The third category is parameter tuning by independent users. The
three sensitivity parameters acknowledged in
Section~\ref{sec:methods-reproducibility} as not formally
optimised, namely the CNN blending weight $\alpha = 0.70$, the GPD
threshold percentile at the $80$th, and the density saturation
count at $2$, are exposed in the configuration file and can be
swept by a user who has a specific application target in mind.
A user can run the same sensitivity sweep on their own region
with the notebooks provided. This shifts the work of optimisation
from a
single best-guess setting in the original paper to a tunable
artefact whose default values are documented and whose tuning
procedure is reproducible.

The tuning is not free: each
parameter sweep multiplies the runtime by the number of values
in the sweep grid. A full three-parameter sweep at five values
per parameter is $125$ end-to-end runs, which is feasible on a
small compute cluster but not on a single workstation. The
open release lowers the entry cost of tuning from re-implementing
the pipeline to changing a configuration line.

A tangible example of what the open release supports is the
station-density confidence mask itself. It is published as a
NetCDF file alongside the corrected product and the source code,
so an independent user can inspect the spatial distribution of
correction confidence directly and verify that the CNN
contribution is suppressed in their region of interest if the
local gauge density is below the saturation threshold.
Figure~\ref{fig:confmask} shows the confidence mask and its
operational consequence on Composite Quality Index (CQI) improvement.
The CQI gain is concentrated in the highest-confidence regions, and
the lowest confidence quintile shows a near-zero median improvement,
by design, because the framework reverts to LSEQM where the reference
cannot support CNN refinement.

\begin{figure}[H]
\centering
\includegraphics[width=\figFiveThreeWidth\textwidth]{fig_thesis_05_confidence_mask.png}
\caption[Station-density confidence mask]{Station-density confidence mask\label{fig:confmask}}
\end{figure}

The fourth category combines comparison studies and training
use. The notebook structure separates correction from
verification, so the four-pillar verification report can be run
on any pair of satellite and reference fields without re-running
the correction. This lets independent groups compare IMERG Final
Run against IMERG Late Run, or a bias-corrected output from a
different pipeline against the same BMKG validation set, against
a numerically comparable baseline. The same structure, combined
with the free-tier Colab runtime in the order of minutes per
dekad, supports use as a graduate-level training resource on
satellite precipitation bias correction
\citep{ref:perkel2018, ref:wagemann2022}. This is not the central
scientific contribution of the present work, but it is a
substantive secondary benefit of the open release.

The reproducibility contribution carries its own limitations,
which Section~\ref{sec:discussion-limitations} catalogues. The
pinned dependency versions will drift, the Colab environment will
evolve, and the input data archives at the National Aeronautics
and Space Administration (NASA) Goddard Earth Sciences Data and
Information Services Center (GES DISC) and National Oceanic and
Atmospheric Administration (NOAA) CPC have their own retention and
re-processing policies that the
present work cannot control. The reproducibility claim is
therefore scoped: it holds against the specific versions and
archive states recorded at the time of release, and it requires
the dependency snapshot to be maintained for the verification
exercise to remain reproducible in the long term. The discussion
returns to this scope in Section~\ref{sec:discussion-limitations}.


% ============================================================
\section{Paths Forward}\label{sec:discussion-paths}
% ============================================================

Raising the temporal-skill ceiling beyond $r \approx 0.35$
requires methods outside the marginal correction family. This
section surveys four paths and assesses their feasibility for the
Indonesian domain, the data products available, and the
infrastructure costs they impose on operational users. The order
reflects increasing methodological distance from the present
pipeline. Table~\ref{tab:paths-forward} summarises the four
paths against three criteria: the methodological distance from
the LSEQM+DL pipeline, the infrastructure cost they impose, and
their effect on the temporal-skill ceiling for the Maritime
Continent setting.

\begin{table}[H]
\centering
\caption[Four paths for raising the temporal-skill ceiling beyond $r \approx 0.35$]{Four paths for raising the temporal-skill ceiling
beyond $r \approx 0.35$\label{tab:paths-forward}}
\begingroup\small
% Use tabular* so the table fills \textwidth exactly: keep the original
% 3.0/2.0/2.5/4.5 cm p-column widths (sum 12 cm) and let
% \extracolsep{\fill} absorb the ~2 cm of slack into the inter-column
% gaps. Previous attempt scaled the p widths to 14 cm but forgot the
% three 12 pt inter-column paddings (~1.3 cm), pushing the table past
% the right margin.
\begin{tabular*}{\textwidth}{@{\extracolsep{\fill}}p{3.0cm}p{2.0cm}p{2.5cm}p{4.5cm}@{}}
\toprule
Path & Methodological distance & Infrastructure cost & Ceiling effect \\
\midrule
Multivariate QM &
Moderate (replace univariate marginal with joint mapping) &
Larger calibration samples; tens of times more compute &
Open question; depends on joint structure of the reference \\
\addlinespace
Sub-daily disaggregation &
Moderate (add a step after daily correction) &
Half-hourly archive + geostationary IR &
Moderate; recovers within-day timing only \\
\addlinespace
Stochastic disaggregation &
Large (deterministic $\to$ ensemble) &
Ensemble verification stack on the downstream side &
Reframes timing as uncertainty rather than removing it \\
\addlinespace
Local-day re-aggregation (with Himawari-9) &
Small (re-aggregate input before pipeline) &
Re-run of pipeline on re-aggregated archive &
Targets the UTC-vs-local-day calendar artefact; most actionable for UTC+7--9 \\
\bottomrule
\end{tabular*}
\endgroup
\end{table}

The first path is multivariate quantile mapping. The univariate
methods used throughout the present pipeline correct the marginal
distribution at each pixel and each dekad independently.
Multivariate methods correct the joint distribution across pixels,
across pairs of variables, or across time lags, and can in
principle restore some of the day-by-day pairing that the
univariate methods leave untouched. The rank preservation
argument of Section~\ref{sec:discussion-bound} is specifically a
property of marginal mapping, so a multivariate method is not
bound by it. The trade-off is that multivariate methods require
substantially larger calibration samples and substantially more
computation, and the quality of the joint structure they recover
is bounded by the joint structure of the calibration reference.
For a regional pipeline calibrated against a $0.5^{\circ}$ gauge
analysis with $172$ validation stations, the joint structure
recoverable is limited, and the gain in temporal skill that a
multivariate method can deliver in this setting is an open
question that the present work does not resolve.

The second path is sub-daily disaggregation. The IMERG product
family is natively half-hourly, and the daily aggregation step
applied in the present pipeline discards the within-day timing
information. Sub-daily disaggregation, in which the corrected
daily total is redistributed across the day using a
cloud-tracking algorithm against half-hourly imagery, recovers
the within-day timing at the cost of additional dependencies on
geostationary infrared and visible channels. For the Maritime
Continent, the Himawari-9 geostationary platform~\citep{ref:bessho2016} provides
$10$-minute cadence at $2$~km resolution and is the natural
candidate for the cloud-tracking input. This path retains the
distributional correction at the daily scale and adds a
within-day timing signal on top, so the bias-free distributional
properties of the present pipeline are preserved while the
day-by-day timing is recovered.

The third path is stochastic disaggregation. Rather than
producing a single corrected daily field, a stochastic
disaggregation step samples from an ensemble of plausible daily
timings consistent with the corrected distributional properties.
The verification posture shifts from deterministic point-by-point
scoring to ensemble metrics: the Continuous Ranked Probability
Score for ensembles, the rank histogram for calibration, and
threshold reliability diagrams for event detection~\citep{ref:wilks2020}. The path is
attractive because it acknowledges the timing uncertainty
explicitly rather than burying it in a flat $r$ value. The cost
is that downstream operational users need to ingest ensemble
inputs, which not all hydrological models or drought-monitoring
applications are equipped to handle. The path is therefore most
useful for applications that are already ensemble-aware.

The fourth path, and the most operationally actionable for users
in the UTC+7 to UTC+9 band, is morning-observation re-aggregation
of the daily IMERG-L archive. The calendar-window mismatch
diagnosed in Section~\ref{sec:discussion-subdaily} places the
optimal alignment at $h \approx -23$~h relative to UTC midnight,
matching the date label the BMKG validation network assigns under
its morning-observation convention \citep{ref:bmkg2006sk32}. The
CPC-UNI calibration target needs no such shift: it is already
dated to the UTC day and peaks at $h \approx 0$
\citep{ref:cpcuni_readme}. Re-aggregating the native half-hourly
IMERG fields \citep{ref:imerghhl07} to a 24-hour window ending at
$07{:}00$ local time on the recorded date, before applying the
LSEQM+DL correction, would remove the mismatch component without
changing the correction methodology.

The expected operational
effect, from the diagnostic in
Section~\ref{sec:discussion-subdaily}, is a lift in the GPM-era
daily Pearson $r$ against BMKG, pooled over all station-days, from
approximately $0.20$ at the
UTC-day window to approximately $0.57$ at the matched window. The
four-pillar corrections apply unchanged, and their value,
distribution, extreme, and detection scores are recomputed on the
re-aggregated daily totals and expected to remain close. The cost is a re-run
of the entire pipeline against a re-aggregated input archive,
which is computationally modest but archive-intensive. This is
the path most directly enabled by the parallel reproducibility
contribution: the configurable daily-aggregation step is exposed
in the pipeline, the re-aggregated archive can be substituted at
the input layer, and the four-pillar verification report on the
re-aggregated pipeline against the same BMKG validation set is
reproducible by any independent user.

A common feature of all four paths is that they leave the
LSEQM+DL pipeline as constructed in Chapter~\ref{chap:methods}
intact for the well-served application class identified in
Section~\ref{sec:discussion-ops}. Users who do not need
day-specific timing receive the corrected daily product as is;
users who need timing receive an upstream re-aggregation, an
add-on disaggregation, or a shift to an ensemble representation.
The marginal correction component remains the right tool for the
distribution-alignment work, and the four paths above add
capability rather than replace existing capability.


% ============================================================
\section{Limitations and Constraints}\label{sec:discussion-limitations}
% ============================================================

The present work has six limitations that bound the strength of
its conclusions. They are catalogued here so that the
operational guidance in Section~\ref{sec:discussion-ops} and the
reproducibility claim in
Section~\ref{sec:discussion-reproducibility} can be read at the
right level of confidence.

\textbf{Sensitivity parameters not formally optimised.} Three
parameters in the pipeline carry the not-formally-optimised
status declared in Section~\ref{sec:methods-reproducibility}: the
CNN blending weight $\alpha = 0.70$, the GPD threshold at the
$80$th percentile of wet-day exceedances, and the density
saturation count at $2$ (lowered from an initial $3$ after
qualitative inspection showed that $3$ suppressed the CNN almost
everywhere outside Java). Each was chosen by domain reasoning,
each is exposed in the configuration file, and each is acknowledged
as a sensitivity parameter rather than an optimum. The headline
numbers in Chapter~\ref{chap:results} are conditional on these
choices. A one-at-a-time sweep over five values per parameter on
the Bali reproducibility subdomain
(Section~\ref{sec:results-sensitivity}) confirms that each
parameter response is monotone and locally smooth and that the
four-pillar pattern does not reverse across the explored ranges; a
full Cartesian grid search with a strict temporal hold-out remains
future work. Users who care about the precise headline value
should re-run the pipeline with their own preferred settings;
users who care about the four-pillar pattern can rely on its
stability across the Bali sweep, which is a single-subdomain test.

\textbf{CPC-UNI reference quality bound.} The framework
calibrates against CPC-UNI, whose known undercatch of heavy
rainfall and $0.5^{\circ}$ interpolation smoothing bound the
reference quality. Section~\ref{sec:methods-reference-quality}
traces both to their source: the Indonesian input reaches CPC
through the routine GTS exchange rather than the full national
network, and that input is drawn from an archive missing a third
of its station-days, so the reference is well constrained where
the contributing network is dense and progressively less so where
it thins. The
corrected satellite product cannot be
better than the reference it is calibrated against, except where
the BMKG out-of-sample validation reveals that the corrected
product matches the BMKG distribution more closely than CPC-UNI
does, as discussed in Section~\ref{sec:results-headline}. The
strongest claim the present work supports is that the corrected
product is bias-free at the four pillars when scored against an
independent gauge network, conditional on the CPC-UNI reference
defining the correction target. Replacement of CPC-UNI by a
higher-resolution or denser-gauge analysis remains an extension
within the same pipeline architecture.

\textbf{CNN training-data limits at the network edges.} The CNN
refinement stage uses dekadal pooling of approximately $250$
paired daily observations per pixel over $25$ years. In regions
of very sparse gauge coverage, the station-density confidence
mask formula in Section~\ref{sec:methods-blending} drives
$\alpha_{\text{eff}}$ toward $1$, which restricts the CNN
contribution to areas where the calibration sample is large
enough to support the residual learning target. This is graceful
degradation rather than failure: in low-confidence regions the
pipeline collapses to LSEQM, and the four-pillar performance in
those regions is the LSEQM column of Table~\ref{tab:headline}
rather than the LSEQM+DL column. The implication for users in
Papua and parts of Maluku is that the headline LSEQM+DL numbers
do not apply uniformly across the archipelago; in those regions
the LSEQM column of Table~\ref{tab:headline} is the relevant
performance baseline.

\textbf{Conditional bias not removable by univariate quantile
mapping.} The Pearson $r$ ceiling discussed in
Section~\ref{sec:discussion-bound} is one expression of a more
general limit: conditional biases, in which the satellite error
depends on the precipitation amount in a way that the univariate
mapping cannot capture, are not removable by the present
pipeline. The four paths in Section~\ref{sec:discussion-paths}
all address this limit at different points in the pipeline. The
present work does not claim to remove conditional bias; it
claims to remove the marginal bias and to acknowledge the
conditional component as the structural ceiling.

\textbf{Calendar-window mismatch between IMERG-L and the gauge
reference family.} The daily IMERG-L archive is constructed on
the UTC calendar day. CPC-UNI shares that date label, closing each
total at $00{:}00$~UTC \citep{ref:cpcuni_readme}, whereas the BMKG
network dates its totals to the local day ending at the morning
observation \citep{ref:bmkg2006sk32}, a mismatch of approximately
$23$~hours per the diagnostic in
Section~\ref{sec:discussion-subdaily}. The mismatch therefore
arises only at the validation end, not at the calibration end. The four-pillar results
(value, distribution, extreme, detection) are unaffected, being
functions of the marginal distribution and of categorical events
that are invariant to the 24-hour window choice; only Pearson~$r$
against the BMKG validation gauges is suppressed by the label
offset, and it is recovered by re-aggregating the satellite to the
local-observation day, identified as the most operationally
actionable path forward (Section~\ref{sec:discussion-paths}).

\textbf{Reproducibility scope is version-bounded.} The
dependencies pinned at release time will drift as the Python
ecosystem evolves and as cloud-platform defaults change. The
reproducibility claim is therefore scoped to the dependency
snapshot recorded in the requirements file at release. A long-term
reproducibility commitment requires either a containerised
distribution (which the present release does not provide) or an
ongoing maintenance commitment to track upstream changes (which
the present author cannot offer beyond the thesis defence
horizon). The verification exercises that
Section~\ref{sec:discussion-reproducibility} discusses should
therefore be undertaken with the recorded dependency snapshot,
not with a current environment, until a containerised release is
prepared.

The combined effect of these six limitations leaves the central
claim of the chapter intact: marginal bias correction moves the
four pillars cleanly toward the gauge target but does not improve
the temporal-skill diagnostics beyond single-pixel sampling noise,
because rank preservation and variance inflation without skill hold
Pearson $r$ near the timing quality of the underlying satellite
retrieval. The four-pillar
improvement is the contribution that the framework delivers; the
flat $r$ is the structural property of the method family rather
than a failure of the implementation. The reproducibility
contribution lowers the cost of testing this proposition on other
products, other references, and other regions; the transferability
of the proposition itself is an empirical question that the open
release lets the field answer one re-application at a time.
Chapter~\ref{chap:conclusion} draws these threads together, states
what the framework does and does not deliver against the objectives
and hypotheses of Chapter~\ref{chap:introduction}, and sets out the
recommendations that follow for users,
for researchers extending the method, and for the institutions that
maintain the underlying archives.
